\documentclass[11pt]{article}
\usepackage[a4paper,margin=25mm]{geometry}
\usepackage{amsmath,amssymb,amsthm,booktabs,xurl,hyperref}
\hypersetup{hidelinks}
\newtheorem{theorem}{Theorem}
\newtheorem{lemma}[theorem]{Lemma}
\newtheorem{corollary}[theorem]{Corollary}
\newcommand{\Q}{\mathbb Q}
\newcommand{\Tr}{\operatorname{Tr}}
\newcommand{\Nm}{\operatorname{N}}
\title{A rational trace in place of each private detector product}
\author{Z diophantine in Q project}
\date{23 November 2026}
\begin{document}
\maketitle
\begin{abstract}
We prove an algebraic compression for any number of quadratic-field
product and norm--trace constraints with a common factor. A single
rational trace replaces each private field product; one common field
is added. For twelve blocks this saves ten rational coordinates and
ten quadratic equations. Applied after the common contact, scalar
point, and scalar norm reductions, it gives complete quadratic tables
in 1112 and 1292 variables. The separate Pell reduction then gives
1103 and 1283 variables. The proof includes both coordinate maps and
all exceptional fibers. The arithmetic interfaces and the unresolved
source universality assertion of the preceding work are retained.
\end{abstract}

\section{An exact algebraic compression}

Let \(K/F\) be a quadratic field extension of characteristic different
from two, with conjugation, trace, and norm denoted by
\(\overline{\phantom{x}}\), \(\Tr\), and \(\Nm\). In the application
\(F=\Q\) and \(K=\Q(\sqrt2)\). Let \(u,h\in F\) be nonzero
common parameters or already guarded coordinates. For
\(i=1,\ldots,r\), suppose \(a_i,W_i,B\in K\) and \(D_i\in F\)
satisfy
\begin{equation}\label{eq:old}
       a_i B=uW_i,\qquad
       (D_i+h)\Tr(W_i)=\Nm(W_i)+h^2.
\end{equation}
The field \(B\) is common. Assume that the private coordinates
\(W_i\) occur nowhere else in the system. All other occurrences
of \(a_i,B,D_i,u,h\) may be retained without change.

\begin{theorem}\label{thm:compression}
For every \(r\ge1\), system \eqref{eq:old} has exactly the same
projection to its other coordinates as the system with one common
field \(v\in K\), rational coordinates \(t_i\in F\), and equations
\begin{equation}\label{eq:new}
 \overline B v=hu,\qquad
 h t_i=\Tr(v a_i),\qquad
 (D_i+h)t_i=\Nm(a_i)+\Nm(v).
\end{equation}
All equations are quadratic in their coordinates. Over a quadratic
extension, this replaces \(2r\) rational product coordinates by
\(r+2\) rational coordinates and replaces \(3r\) rational equations
by \(2r+2\). Thus it saves \(r-2\) coordinates and equations.
\end{theorem}

\begin{proof}
If an old solution had \(W_i=0\), its second equation would give
\(h^2=0\), impossible. Hence \(W_i\ne0\), and the first equation
forces \(B\ne0\). Because \(K\) is a field, \(\Nm(B)\ne0\).
Define
\begin{equation}\label{eq:forward}
             v=\frac{hu}{\overline B},\qquad
             t_i=\frac{\Tr(v a_i)}{h}.
\end{equation}
Then \(\overline Bv=hu\). The identities
\[
 \frac1{\overline B}=\frac{B}{\Nm(B)},\qquad
 t_i=\frac{u\Tr(a_iB)}{\Nm(B)},\qquad
 \Nm(v)=\frac{h^2u^2}{\Nm(B)}
\]
give the exact residual identity
\begin{align}\label{eq:identity}
 &(D_i+h)t_i-\Nm(a_i)-\Nm(v)\nonumber\\
 &\qquad=\frac{u^2}{\Nm(B)}
 \bigl((D_i+h)\Tr(W_i)-\Nm(W_i)-h^2\bigr),
 \qquad W_i=\frac{a_iB}{u}.
\end{align}
This proves all equations of the new system.

Conversely, start with a new solution. The equation
\(\overline Bv=hu\ne0\) forces \(B,v\ne0\). Define
\[
                            W_i=\frac{a_iB}{u}.
\]
The common equation and the trace equation imply the identities
in \eqref{eq:forward}. The nonzero multiplier
\(u^2/\Nm(B)\) in \eqref{eq:identity} then proves the old
norm--trace equation. This recovers every removed equation.

The maps divide only by \(u,h,B\), whose nonvanishing has just
been proved or was already guarded. In particular no division by
\(t_i\), \(D_i+h\), or a private component of \(a_i\) is made.
The equivalence therefore concerns the whole solution set, not a
generic open subset chosen after discarding exceptional fibers.
Finally, one field equation has two rational components, while a
trace or norm equation is a single rational equation. This proves
the stated counts.
\end{proof}

\section{The ordinary quadratic component equations}

Write
\[
 B=b_0+b_1\sqrt2,\qquad v=v_0+v_1\sqrt2,
 \qquad a_i=a_{i0}+a_{i1}\sqrt2.
\]
The common field equation in \eqref{eq:new} is precisely
\begin{equation}\label{eq:components-common}
                  b_0v_0-2b_1v_1=hu,\qquad
                  b_0v_1-b_1v_0=0.
\end{equation}
The two equations for block \(i\) are
\begin{equation}\label{eq:components-private}
\begin{aligned}
 h t_i&=2v_0a_{i0}+4v_1a_{i1},\\
 (D_i+h)t_i&=a_{i0}^2-2a_{i1}^2+v_0^2-2v_1^2.
\end{aligned}
\end{equation}
No separate coordinate for \(\Nm(v)\) is needed: it is already
a quadratic expression and appears with coefficient one. Similarly,
\(hu\) is quadratic in the existing coordinates. These observations
are essential for obtaining the full saving \(r-2\).

\section{Application to the complete combined systems}

The common-contact construction stores nonzero rational coordinates
\(c,c^{-1},c^{-2}\) with the literal quadratic guards
\[
                    c\,c^{-1}=1,\qquad c^{-2}=(c^{-1})^2.
\]
In both coefficient families put \(u=c^{-2}\). In the general
family take \(h=c^{-2}\); in the specified family take
\(h=c^{-1}\). Both are already stored and nonzero.

The product of the three common detector factors is the common
field \(B\). The remaining factor is the private field \(a_i\).
The retained final multiplication rows and norm-error row are
exactly \eqref{eq:old}. In the specified table the literal constant
in the norm row is \(c^{-2}\), equal to \(h^2\) by the retained
unit guard. An exhaustive occurrence check confirms that the two
components of each \(W_i\) occur in exactly those three rows.
Thus Theorem~\ref{thm:compression} applies without changing any
contact gate, arithmetic estimate, scalar block, or input chart.

The default complete predecessors are the scalar-norm tables
\cite{norm}, which already include the common-contact and
scalar-point reductions \cite{global,point}. There are twelve
blocks. The finite counts are
\begin{center}
\begin{tabular}{lrrrr}
\toprule
Interface & Old variables & New variables & Old equations & New equations\\
\midrule
General &1122&1112&966&956\\
Specified &1302&1292&1122&1112\\
\bottomrule
\end{tabular}
\end{center}

The nine-coordinate, three-equation Pell reduction
\cite{pell} affects disjoint coordinates and rows. It therefore
commutes with this change and gives the further complete counts
\[
 \begin{array}{c|cc}
 &\text{rational variables}&\text{quadratic equations}\\\hline
 \text{general}&1103&953\\
 \text{specified}&1283&1109.
 \end{array}
\]
The displayed tables before the Pell step and the combined tables
after that step have distinct files; the Pell saving must be counted
only once.

The emended Jones source still carries the unresolved universality
hypothesis \(U_{\rm em}\). Its frozen parameter involving \(n^2\)
has not been altered here. The arithmetic interfaces of the
predecessor remain hypotheses as stated in the cited project papers.
This exact equivalence introduces no \(abc\) assumption, but it
does not remove an existing unproved hypothesis or establish the
claimed 57-variable universal source. In particular, the counts
above are certified complete quadratic checkpoints, with exactly
the conditional mathematical interpretation of their predecessors.

\section{Reproducibility and verification}

The implementation
\nolinkurl{undecidability/code/build_detector_trace.py}
reads a complete predecessor, removes only the audited product
coordinates and rows, and writes every new ordinary quadratic
residual. Its optional baseline-directory and baseline-prefix
arguments permit composition with the independent source-core
changes. The default output is stored in
\nolinkurl{undecidability/coefficients/jones1980_detector_trace/}.

The independent checker
\nolinkurl{undecidability/verify_detector_trace.py}
reconstructs every old product and norm row from the field formulas,
checks the literal unit equations, verifies the complete occurrence
lists, and reconstructs every row of the new table. It also checks
\eqref{eq:identity} as an exact rational-function identity in
independent indeterminates, the coordinate census, degree at most
two, integer coefficient domain, complete coordinate support,
canonical rebuilt bytes, and the predecessor and result hashes.
The table certificates explicitly retain the unverified source
status. No numerical testing is used to replace a proof of
equivalence or a coefficient identity.

Over \(\Q\), the sum of the squares of the listed residuals
vanishes exactly when every residual vanishes. The unchanged outer
system contains nonzero quadratic leading terms, so that sum has
degree exactly four. The coefficient files specify this quartic
without requiring an expanded list of its monomials.

\begin{thebibliography}{9}
\bibitem{global} Z diophantine in Q project,
\emph{One contact gauge in a combined quadratic system},
\nolinkurl{undecidability/papers/global_gauge_quartic_checkpoint.tex}.
\bibitem{point} Z diophantine in Q project,
\emph{Removing private scalar normalizers from a common point chart},
\nolinkurl{undecidability/papers/scalar_rechart_quartic_checkpoint.tex}.
\bibitem{norm} Z diophantine in Q project,
\emph{Scalar norm compression},
\nolinkurl{undecidability/papers/scalar_norm_compression.tex}.
\bibitem{pell} Z diophantine in Q project,
\emph{Pell index optimization},
\nolinkurl{undecidability/papers/pell_index_optimization.tex}.
\end{thebibliography}
\end{document}
